\begin{tikzpicture}[
  >=Latex,
  line cap=round,
  line join=round,
  font=\small,
  gravedad/.style={->, orange!85!black, line width=1pt},
  desplazamiento/.style={->, violet!80!black, line width=1.2pt},
  inicial/.style={densely dashed, gray!65, line width=1pt}
]
% Panel (a): corte por el centro de una esfera.
\node[font=\small\bfseries] at (0,2.65) {(a) Contracción gravitacional};
\draw[inicial] (0,0) circle (2.0);
\draw[blue!65!black,line width=1.3pt] (0,0) circle (1.52);
\fill[gray!35] (0,0) circle (1.52);
\draw[blue!65!black,line width=1.3pt] (0,0) circle (1.52);
\foreach \angulo in {0,45,90,135,180,225,270,315} {
  \draw[gravedad] ({1.82*cos(\angulo)},{1.82*sin(\angulo)})
    -- ({0.88*cos(\angulo)},{0.88*sin(\angulo)});
  \draw[desplazamiento] ({1.52*cos(\angulo)},{1.52*sin(\angulo)})
    -- ({1.23*cos(\angulo)},{1.23*sin(\angulo)});
}
\fill[black] (0,0) circle (1.5pt) node[below right=2pt] {centro};
\draw[->,gray!70] (0,0) -- ({2*cos(35)},{2*sin(35)}) node[pos=0.72,above] {$R$};
\node[gray!70,anchor=west] at (-1.95,-2.32) {radio inicial (discontinuo)};
\node[blue!65!black,anchor=west] at (-1.95,-2.62) {radio deformado (exagerado)};

% Panel (b): campo gravitacional de densidad uniforme.
\begin{scope}[xshift=6.8cm]
  \node[font=\small\bfseries] at (1.65,2.65) {(b) Magnitud de $g(r)$};
  \draw[->,gray!75] (-0.15,0) -- (3.8,0) node[right] {$r$};
  \draw[->,gray!75] (0,-0.15) -- (0,2.35) node[above] {$g$};
  \draw[blue!65!black,line width=1.3pt] (0,0) -- (3.1,1.85);
  \fill[blue!65!black] (0,0) circle (2pt) (3.1,1.85) circle (2pt);
  \draw[gray!55,densely dotted] (3.1,0) -- (3.1,1.85) -- (0,1.85);
  \node[below] at (3.1,0) {$R$};
  \node[left] at (0,1.85) {$g_R$};
  \node[below left] at (0,0) {$0$};
  \node[font=\footnotesize,align=center] at (1.6,-0.72)
    {$g(r)=g_R\,r/R$ para densidad uniforme};
\end{scope}
\end{tikzpicture}